% ============================================================
%  One file per finding. Copy this file to add a new finding and
%  update the \input list in main.tex.
% ============================================================
\newpage
\subsection{F-01 --- \fillme{Short, specific finding title, e.g. "Unauthenticated SQL Injection in the /search Endpoint"}}
\label{sec:f01}

\begin{center}
\begin{tabularx}{\linewidth}{@{}p{3.8cm} X p{3.8cm} X@{}}
\toprule
\rowcolor{TableHead}
\multicolumn{4}{l}{\textcolor{white}{\bfseries\sffamily Finding Metadata}} \\
\midrule
\textbf{ID}            & F-01                                & \textbf{Severity}             & \critical \\
\rowcolor{TableAlt}
\textbf{Affected Endpoint} & \ep{/api/v1/search?q=} & \textbf{Exploitability} & \fillme{Trivial / Moderate / Difficult} \\
\textbf{CVSS Score}    & \fillme{9.8 (Critical)}             & \textbf{Authentication}       & \fillme{None / Authenticated} \\
\rowcolor{TableAlt}
\textbf{Category}      & \fillme{SQL Injection (CWE-89)}      & \textbf{Remediation Priority} & \fillme{Immediate} \\
\bottomrule
\end{tabularx}
\end{center}

\subsubsection{Description}

% <EDIT> Explain, in plain language, what the vulnerable component does and how the
% attacker-controlled input reaches a dangerous sink.
\fillme{Describe the vulnerable endpoint/parameter, what it does normally, and how untrusted
input reaches the sink without proper validation or parameterisation.}

\subsubsection{Steps to Reproduce}

% <EDIT> Number the steps you actually took. Reference a screenshot per step where useful;
% drop a PNG under images/<target>/ and reference it exactly as below.
\begin{enumerate}
  \item \fillme{Step 1 -- e.g. send a baseline request and note the normal response.}
\begin{figure}[H]
  \centering
  % \includegraphics[width=0.85\textwidth]{example.com/01-baseline-response.png}
  \fbox{\parbox{0.85\textwidth}{\centering\vspace{2cm}\fillme{screenshot placeholder --
    replace with \textbackslash includegraphics once you add the image under images/}\vspace{2cm}}}
  \caption{\fillme{Baseline: normal response from the endpoint}}
\end{figure}
  \item \fillme{Step 2 -- e.g. supply the payload and note the changed behaviour.}
  \item \fillme{Step 3 -- e.g. confirm and characterise with tooling.}
\end{enumerate}

\subsubsection{Impact}

\begin{criticalbox}
\fillme{One tight paragraph: what an attacker can concretely do as a result of this finding,
stated in business terms (data exposed, systems reachable, accounts takeover-able), and how it
relates to any other finding it chains with.}
\end{criticalbox}

\subsubsection{Remediation}

\begin{itemize}
  \item \fillme{The specific code/config fix, e.g. "use parameterised queries / an ORM everywhere this pattern occurs".}
  \item \fillme{A defence-in-depth control, e.g. least-privilege DB account, WAF rule as an interim mitigation.}
\end{itemize}
